\section{STERA 30M}
\label{sec:architettura-stera30m}

STERA 30M è un Transformer decoder-only a sei strati, addestrato da zero per \num{122000} aggiornamenti. Ha dimensione nascosta 768, 16 teste di attenzione da 48 componenti e un livello feed-forward da 3.072 componenti con GELU. Alterna tre strati di attenzione causale sull'intero prefisso e tre con finestra locale di 256 posizioni. Il training usa finestre fino a 512 token. La tabella delle posizioni assolute contiene 2.048 righe allenabili.

\begin{figure}[htbp]
\centering
\resizebox{.92\linewidth}{!}{\begin{tikzpicture}[
  font=\small, >={Latex[length=2mm]},
  box/.style={draw=steraGrigio, rounded corners=2pt, text width=60mm, minimum height=7mm, align=center, fill=white},
  rm/.style={box, fill=steraBlu!12, draw=steraBlu},
  op/.style={circle, draw=steraGrigio, inner sep=1pt, font=\footnotesize},
  lab/.style={font=\footnotesize, text=steraGrigio, align=left}]
\node[box] (in) {Token $u_t$ e posizione $t$};
\node[rm, above=5mm of in] (emb) {Embedding: $x_t = E_{u_t,:} + P_{t,:}$};
\node[box, above=8mm of emb] (ln1) {LayerNorm};
\node[rm, above=3mm of ln1] (att) {Attenzione causale, 16 teste\\[-1pt]{\footnotesize $W_Q, W_K, W_V, W_O$ RowMask23}};
\node[op, above=3mm of att] (s1) {$+$};
\node[box, above=3mm of s1] (ln2) {LayerNorm};
\node[rm, above=3mm of ln2] (ff) {Feed-forward con GELU, $768\to3072\to768$\\[-1pt]{\footnotesize $W_1, W_2$ RowMask23}};
\node[op, above=3mm of ff] (s2) {$+$};
\node[box, above=8mm of s2] (lnf) {LayerNorm finale};
\node[rm, above=5mm of lnf] (out) {Logit: $z_t = E\,h_t$ \quad {\footnotesize (50.257 candidati)}};
\node[box, above=5mm of out] (p) {$p(u_{t+1}\mid u_{\le t}) = \mathrm{softmax}(z_t)$};
\draw[->] (in) -- (emb); \draw[->] (emb) -- (ln1); \draw[->] (ln1) -- (att); \draw[->] (att) -- (s1);
\draw[->] (s1) -- (ln2); \draw[->] (ln2) -- (ff); \draw[->] (ff) -- (s2); \draw[->] (s2) -- (lnf);
\draw[->] (lnf) -- (out); \draw[->] (out) -- (p);
\coordinate (r0) at ($(emb.north)!.5!(ln1.south)$);
\draw[->] (r0) -- ++(-38mm,0) |- (s1.west);
\coordinate (r1) at ($(s1.north)!.5!(ln2.south)$);
\draw[->] (r1) -- ++(-38mm,0) |- (s2.west);
\begin{scope}[on background layer]
\node[draw=steraGrigio, dashed, rounded corners=4pt, fit=(ln1)(att)(s1)(ln2)(ff)(s2), inner xsep=10mm, inner ysep=3mm] (blk) {};
\end{scope}
\node[lab, anchor=north west] at ([xshift=2mm,yshift=-2mm]blk.north east) {$\times 6$ strati};
\node[lab, anchor=north west, text width=27mm] at ([xshift=2mm,yshift=-8mm]blk.north east) {strati 1, 3, 5: attenzione su tutto il prefisso\\ strati 2, 4, 6: finestra locale di 256 posizioni};
\draw[->, steraBlu, dashed] (emb.west) -- ++(-20mm,0) |- node[pos=.25, left, font=\footnotesize, text=steraBlu]{stessa $E$} (out.west);
\end{tikzpicture}
}
\caption{Architettura di STERA 30M, dal token in ingresso alla distribuzione sul token successivo. In azzurro le matrici rappresentate con RowMask23; la matrice $E$ è condivisa fra embedding e testa di uscita.}
\label{fig:stera30m-flusso}
\end{figure}

\subsection{Dal token ai punteggi del vocabolario}

Il tokenizzatore associa a ogni token un indice $u_t$ fra 0 e 50.256. Il modello legge la riga $E_{u_t,:}$ e vi somma il vettore di posizione $P_{t,:}$. Il vettore risultante ha 768 componenti e attraversa i sei strati. Dopo la normalizzazione finale, chiamiamo $h_t$ la rappresentazione contestuale alla posizione $t$. Il punteggio del candidato $v$ è $\ell_{t,v}=E_{v,:}h_t$. La softmax sui \num{50257} logits restituisce le probabilità del token successivo.

In questo percorso lo stesso elemento di $E$ può influire sia sul vettore che entra nel Transformer sia sul punteggio di un candidato in uscita. La matrice fisica ha \num{50304} righe per allinearsi ai blocchi da 128. Le 47 righe aggiuntive sono escluse dal vocabolario e dal calcolo delle probabilità.

\subsection{Come si costruiscono le matrici}

Ogni matrice principale contiene un mosaico di blocchi RowMask23. Ciascun blocco possiede la propria tabella con 127 valori di base fuori diagonale, 128 valori diagonali e 2.921 correzioni, pari a 23 per ognuna delle righe da 1 a 127. Lo XOR degli indici determina quale valore della tabella usare, mentre la maschera determina dove aggiungere una correzione. La disposizione delle correzioni nel blocco è quella della Figura~\ref{fig:struttura}.

Per una matrice di $m$ righe e $n$ colonne, indichiamo con $R(c,s)$ il mosaico ricostruito da coordinate $c$ e scale $s$, una per blocco. La matrice effettiva è
\begin{equation}
 W = \operatorname{diag}(g_{\mathrm{out}})\,R(c,s)\,
     \operatorname{diag}(g_{\mathrm{in}}) + UV^\top,
 \qquad U \in \mathbb{R}^{m\times96},\quad V \in \mathbb{R}^{n\times96}.
 \label{eq:stera30m-peso}
\end{equation}
I guadagni agiscono per riga e per colonna. Il termine globale $UV^\top$ ha rango al più 96 e può modificare anche posizioni appartenenti a blocchi diversi.

La matrice condivisa $E$ conta $393\cdot6=2.358$ blocchi. Ogni strato ha quattro matrici $768\times768$ da 36 blocchi ciascuna e due proiezioni feed-forward da 144 blocchi ciascuna. Il totale è $1+6\cdot6=37$ matrici e $2358+6(4\cdot36+2\cdot144)=\num{4950}$ blocchi.

\begin{table}[htbp]
\centering
\small
\begin{tabular}{lr}
\toprule
Componente & Parametri allenabili \\
\midrule
Coordinate RowMask23 & \num{15721200} \\
Scale per blocco & \num{4950} \\
Guadagni per riga e colonna & \num{134016} \\
Fattori $U$ e $V$ & \num{12865536} \\
Posizioni assolute & \num{1572864} \\
Parametri di normalizzazione & \num{19968} \\
Bias delle proiezioni & \num{27648} \\
\midrule
Totale & \textbf{\num{30346182}} \\
\bottomrule
\end{tabular}
\caption{Inventario dei parametri di STERA 30M. La matrice condivisa è contata una volta. La rappresentazione espansa contiene \num{82721280} valori, inclusi i parametri già densi.}
\label{tab:stera30m-parametri}
\end{table}

\subsection{Come risale il gradiente}

In una posizione con token bersaglio $y_t$, la perdita scalare è $\mathcal L_t=-\log p_{t,y_t}$. Il gradiente rispetto al logit del candidato $v$ è $d_{t,v}=p_{t,v}-\mathbf{1}_{\{v=y_t\}}$. Con $d_t$ vettore colonna, il contributo della testa al gradiente di $E$ sulle righe del vocabolario è $d_t h_t^\top$, mentre il contributo verso la rappresentazione contestuale è $E_{\mathrm{voc}}^\top d_t$. Qui $E_{\mathrm{voc}}$ indica le sole righe valide di $E$. Attraversando gli strati, il gradiente contribuisce anche alle righe di $E$ usate dai lookup in ingresso. I due contributi si sommano sulla stessa matrice.

\begin{figure}[htbp]
\centering
\begin{tikzpicture}[font=\small, >={Latex[length=2mm]},
  box/.style={draw=steraGrigio, rounded corners=2pt, minimum width=60mm, minimum height=8mm, align=center, fill=white},
  leaf/.style={draw=steraBlu, rounded corners=2pt, fill=steraBlu!10, minimum width=36mm, minimum height=13mm, align=center, font=\footnotesize}]
\node[box] (l) {Perdita $\mathcal{L}_t = -\log p_{t,y_t}$};
\node[box, below=6mm of l] (z) {Logit: $\partial\mathcal{L}/\partial z_{t,v} = p_{t,v} - \mathbf{1}\{v=y_t\}$};
\node[box, below=6mm of z] (h) {Strati Transformer: $G = \partial\mathcal{L}/\partial W$ per ogni proiezione};
\node[box, below=6mm of h] (w) {$W = \mathrm{diag}(g_{\mathrm{out}})\,R(c,s)\,\mathrm{diag}(g_{\mathrm{in}}) + UV^{\top}$};
\node[leaf, below left=9mm and -14mm of w] (d) {Diagonale\\$\partial\mathcal{L}/\partial d_i = \tilde G_{ii}$};
\node[leaf, below=9mm of w] (b) {Base XOR\\$\partial\mathcal{L}/\partial b_s = \sum_{i\oplus j=s}\tilde G_{ij}$};
\node[leaf, below right=9mm and -14mm of w] (e) {Correzioni\\$\partial\mathcal{L}/\partial e_{i,k} = \tilde G_{i,\,i\oplus\sigma(i,k)}$};
\node[font=\footnotesize, text=steraGrigio, align=left, anchor=west] at ([xshift=4mm]w.east) {$\partial\mathcal{L}/\partial U = GV$\\$\partial\mathcal{L}/\partial V = G^{\top}U$};
\draw[->] (l) -- (z); \draw[->] (z) -- (h); \draw[->] (h) -- (w);
\draw[->] (w.south) -- (d.north); \draw[->] (w.south) -- (b.north); \draw[->] (w.south) -- (e.north);
\node[font=\footnotesize, text=steraGrigio, below=3mm of b, align=center] {$\tilde G$: gradiente del blocco moltiplicato per scala e guadagni di riga e colonna};
\end{tikzpicture}

\caption{Retropropagazione dalla perdita scalare alla matrice condivisa e alle coordinate RowMask23.}
\label{fig:stera30m-backward}
\end{figure}

Per una base XOR $b_x$, il gradiente somma i contributi di tutte le caselle del blocco con $i\oplus j=x$. Alla somma concorrono i guadagni delle righe e colonne interessate e la scala del blocco. Una correzione riceve il contributo della propria casella. Se $G=\partial\mathcal L/\partial W$, i gradienti del termine a basso rango sono $GV$ per $U$ e $G^\top U$ per $V$. AdamW aggiorna i tensori allenabili della Tabella~\ref{tab:stera30m-parametri}.

\subsection{Training e motori misurati}

STERA 30M è stato addestrato con un'implementazione PyTorch/CUDA su due GPU NVIDIA P100. Il passo in avanti ricostruisce le matrici effettive. La perdita è calcolata sui token bersaglio, escludendo il riempimento. L'addestramento elabora complessivamente \num{1529514422} token bersaglio.

Per le misure di esecuzione, il percorso denso carica le matrici FP32 già espanse, mentre il percorso compatto applica le proiezioni direttamente alle coordinate tramite kernel Metal. Entrambi eseguono gli stessi pesi. I protocolli di esecuzione sono nel Capitolo~\ref{cap:capitoli/07-sistemi}.
